\documentclass[11pt]{article}
\usepackage[margin=1in]{geometry}
\usepackage{amsmath,amssymb,mathtools,amsthm}
\usepackage{booktabs,array}
\usepackage{enumitem}
\setlist[itemize]{leftmargin=*,itemsep=0.25em,topsep=0.2em}
\setlist[enumerate]{leftmargin=*,itemsep=0.25em,topsep=0.2em}

\newcommand{\R}{\mathbb{R}}
\newcommand{\E}{\mathbb{E}}
\newcommand{\PP}{\mathbb{P}}
\newcommand{\1}{\mathbf{1}}

\title{IE206 Complete Notes (from Provided PDFs)}
\author{}
\date{\today}

\begin{document}
\maketitle
\tableofcontents
\newpage

\section{Coverage Map (All Provided Files)}
\begin{center}
\begin{tabular}{>{\raggedright\arraybackslash}p{0.37\linewidth} >{\raggedright\arraybackslash}p{0.57\linewidth}}
\toprule
Source PDF & Topics included in these notes \\
\midrule
\texttt{IE206\_Lec01.pdf} & Supervised setup, classification vs regression, simple/multiple linear regression, MSE/LSE, correlation, $R^2$, SST/SSM/SSE, polynomial fit \\
\texttt{IE206\_Logistics\_2.pdf} & Model validation perspective, regression diagnostics, outliers/influence, robust losses (L2/L1/Huber/Bisquare), Ridge/LASSO, Bayesian view of regularization, logistic regression and sigmoid/logit \\
\texttt{IE\_206\_ModelEvaluation\_3.pdf} & Hypothesis classes, train/test split, confusion matrix metrics, precision/recall/F1, underfit/overfit, holdout, $k$-fold CV, model selection curve \\
\texttt{IE\_206\_SGD.pdf} & GD/SGD basics, subgradients, hinge subgradient, step-size schedules, strongly convex case, projected SGD, convergence guarantee in convex-Lipschitz bounded setting \\
\texttt{IE\_206\_Perceptron\_4.pdf} & Halfspaces, perceptron update rule, convergence theorem and mistake bound, feature scaling/standardization, feature engineering and polynomial basis \\
\texttt{IE\_206\_SVM.pdf} & Linear separability, margin geometry, hard/soft SVM, slack variables, hinge-loss regularization, support vectors, duality, dual optimization \\
\texttt{IE\_206\_KNN.pdf} & KNN algorithm, distance metrics, curse of dimensionality, Bayes classifier/risk/optimality, Naive Bayes, discriminative vs generative methods \\
\texttt{IE\_206\_MultiClass.pdf} & One-vs-All, All-Pairs, reduction pitfalls, linear multiclass predictors, multiclass SVM, SGD update, ranking intro and losses \\
\texttt{IE\_206\_RandoMForest.pdf} & Decision trees, ID3, entropy/gini gains, threshold splits for real-valued features, random forest construction, boosting, decision stumps, ERM for stumps, AdaBoost and training error bound \\
\texttt{IE206\_PracticeQs\_1.pdf} & Cost-sensitive model choice, imbalanced metrics tradeoff, LOOCV in linear regression, hyperparameter selection, Sherman-Morrison identity, fast LOOCV formula via hat matrix, convexity/nondifferentiability of clipped regression loss \\
\texttt{IE206\_PracticeQs\_2.pdf} & Bayes vs 1-NN error bound, scaling effect in KNN, hard-margin SVM worked geometry, soft-margin extremes ($C=0,\infty$), KKT interpretation ($\alpha_i=C$ case) \\
\bottomrule
\end{tabular}
\end{center}

\section{Linear Regression Notes (\texttt{IE206\_Lec01.pdf})}
\subsection*{Learning setup}
Given i.i.d. samples
\[
S=\{(x^{(i)},y^{(i)})\}_{i=1}^m,
\]
classification has finite/discrete label set $\mathcal Y$, while regression has continuous (or non-countable) labels.

\subsection*{Simple linear regression}
Model:
\[
y=\beta_0+\beta_1x.
\]
Probabilistic form with noise:
\[
y^{(i)}=\beta_0+\beta_1x^{(i)}+\varepsilon_i,
\qquad \E[y^{(i)}\mid x^{(i)}]=\beta_0+\beta_1x^{(i)}.
\]

\subsection*{Goodness of fit and least squares}
Residual $r_i=y^{(i)}-(\beta_0+\beta_1x^{(i)})$. Common losses:
\begin{itemize}
\item absolute error: $|r_i|$
\item squared error: $r_i^2$
\end{itemize}
MSE/LSE objective:
\[
(\hat\beta_0,\hat\beta_1)=\arg\min_{\beta_0,\beta_1}\sum_{i=1}^m r_i^2.
\]
Closed form:
\[
\hat\beta_1=\frac{\sum_i(x_i-\bar x)(y_i-\bar y)}{\sum_i(x_i-\bar x)^2},
\qquad
\hat\beta_0=\bar y-\hat\beta_1\bar x.
\]

\subsection*{Correlation and interpretation}
Sample correlation coefficient:
\[
r=\frac{\sum_i(x_i-\bar x)(y_i-\bar y)}{\sqrt{\sum_i(x_i-\bar x)^2}\sqrt{\sum_i(y_i-\bar y)^2}}.
\]
Interpretation: sign gives direction, magnitude gives linear association strength.

\subsection*{Multiple linear regression}
\[
y^{(i)}=x^{(i)\top}\beta+\varepsilon_i,
\quad
\hat\beta=\arg\min_\beta\sum_{i=1}^m(y^{(i)}-x^{(i)\top}\beta)^2.
\]
Matrix form (standard convention with $X\in\R^{m\times d}$):
\[
\hat\beta=(X^\top X)^{-1}X^\top y
\]
(when $X^\top X$ is invertible).

\subsection*{Coefficient of determination and decomposition}
\[
\text{SST}=\sum_i(y_i-\bar y)^2,
\quad
\text{SSE}=\sum_i(y_i-\hat y_i)^2,
\quad
\text{SSM}=\sum_i(\hat y_i-\bar y)^2.
\]
Identity:
\[
\text{SST}=\text{SSM}+\text{SSE},
\qquad
R^2=1-\frac{\text{SSE}}{\text{SST}}=\frac{\text{SSM}}{\text{SST}}.
\]

\subsection*{Polynomial fit (feature expansion)}
Linear in parameters can represent nonlinear curves by expanding features:
\[
\phi(x)=(1,x,x^2,\dots,x^p),
\quad y\approx w^\top\phi(x).
\]

\subsection*{Important caution}
High training fit does not imply good extrapolation/generalization; always validate on unseen data.

\section{Logistic + Robust Regression Notes (\texttt{IE206\_Logistics\_2.pdf})}
\subsection*{Regression diagnostics and outliers}
Key checks:
\begin{itemize}
\item residual behavior (random vs patterned)
\item outliers (large residual) and influential points (large effect on fitted model)
\item feature redundancy/collinearity
\end{itemize}

\subsection*{Robustness through loss design}
Let $r_i=y_i-x_i^\top\beta$.
\begin{itemize}
\item L2 (squared): $\rho(r)=r^2$ (sensitive to large outliers)
\item LAD/L1: $\rho(r)=|r|$ (more robust; nondifferentiable at 0)
\item Huber:
\[
\rho_k(r)=
\begin{cases}
\frac12r^2,& |r|\le k,\\
k\left(|r|-\frac{k}{2}\right),& |r|>k,
\end{cases}
\]
quadratic near 0 and linear in tails
\item Bisquare/Tukey: strong downweighting of very large residuals
\end{itemize}

\subsection*{Distribution view of loss choice}
\begin{itemize}
\item Gaussian noise $\Rightarrow$ squared loss (MLE)
\item Laplace noise $\Rightarrow$ absolute loss (MLE)
\end{itemize}

\subsection*{Ridge and LASSO}
\textbf{Ridge:}
\[
\hat\beta_{\text{ridge}}
=\arg\min_\beta\sum_{i=1}^m(y_i-x_i^\top\beta)^2+\lambda\sum_{j=1}^d\beta_j^2
\]
\[
\hat\beta_{\text{ridge}}=(X^\top X+\lambda I)^{-1}X^\top y.
\]
Effects: shrinkage, handles collinearity, usually not sparse.

\textbf{LASSO:}
\[
\hat\beta_{\text{lasso}}=
\arg\min_\beta\sum_{i=1}^m(y_i-x_i^\top\beta)^2+\lambda\sum_{j=1}^d|\beta_j|.
\]
Effects: shrinkage + feature selection (sparse coefficients).

\subsection*{Bayesian connection}
MAP estimation with priors:
\begin{itemize}
\item Gaussian prior on $\beta$ $\Rightarrow$ Ridge penalty
\item Laplace prior on $\beta$ $\Rightarrow$ LASSO penalty
\end{itemize}

\subsection*{Logistic regression}
Binary classification with score $\eta=x^\top\beta$:
\[
\PP(Y=1\mid X=x)=\sigma(\eta)=\frac{1}{1+e^{-\eta}}.
\]
Logit transform:
\[
\log\frac{p(x)}{1-p(x)}=x^\top\beta.
\]
Likelihood for data $\{(x_i,y_i)\}$, $y_i\in\{0,1\}$:
\[
L(\beta)=\prod_{i=1}^m p_i^{y_i}(1-p_i)^{1-y_i},
\quad p_i=\sigma(x_i^\top\beta).
\]
Equivalent minimization is negative log-likelihood (logistic loss), solved iteratively (no closed-form like OLS).

\section{Model Evaluation Notes (\texttt{IE\_206\_ModelEvaluation\_3.pdf})}
\subsection*{Hypothesis class}
A hypothesis is a mapping $h:\mathcal X\to\mathcal Y$.
Examples: linear, quadratic, hyperplane, polynomial classes. Learning chooses a good $h$ from $\mathcal H$.

\subsection*{Train-test split}
Typical split: train 70--80\%, test 20--30\%. Test data estimates out-of-sample performance.

\subsection*{Confusion matrix metrics}
With TP, TN, FP, FN:
\[
\text{Accuracy}=\frac{TP+TN}{TP+TN+FP+FN},
\]
\[
\text{Precision}=\frac{TP}{TP+FP},
\qquad
\text{Recall}=\frac{TP}{TP+FN},
\]
\[
F1=\frac{2\cdot\text{Precision}\cdot\text{Recall}}{\text{Precision}+\text{Recall}}.
\]
Use precision when FP cost is high; recall when FN cost is high; F1 for class imbalance/tradeoff.

\subsection*{Underfitting vs overfitting}
\begin{itemize}
\item Underfit: high train and test error (high bias)
\item Proper complexity: low test error
\item Overfit: low train but high test error (high variance)
\end{itemize}

\subsection*{Validation strategies}
\begin{itemize}
\item Holdout: simple, but can waste data and have high variance estimate
\item $k$-fold CV: each sample validated once, more stable estimate
\item Practical default often $k=5$ or $10$ depending on data size/compute
\end{itemize}
Model selection should be based on validation/CV performance, not training error.

\section{SGD Notes (\texttt{IE\_206\_SGD.pdf})}
\subsection*{Gradient descent}
For differentiable $f:\R^d\to\R$:
\[
 w_{t+1}=w_t-\eta\nabla f(w_t).
\]
$\eta$ controls step size/stability.

\subsection*{Subgradient for non-differentiable convex functions}
A vector $g$ is subgradient at $w$ if
\[
f(u)\ge f(w)+\langle g,u-w\rangle\quad\forall u.
\]
Used when gradient does not exist (e.g., hinge, absolute value).

\subsection*{Hinge loss and Lipschitz idea}
For labeled example $(x,y)$:
\[
\ell_{\text{hinge}}(w;(x,y))=\max(0,1-y\langle w,x\rangle).
\]
Subgradient:
\[
\partial \ell=
\begin{cases}
\{0\}, & y\langle w,x\rangle>1,\\
\{-yx\}, & y\langle w,x\rangle<1,\\
\text{conv}\{0,-yx\}, & y\langle w,x\rangle=1.
\end{cases}
\]

\subsection*{Stochastic gradient descent}
Use unbiased stochastic (sub)gradient $v_t$:
\[
\E[v_t\mid w_t]\in\partial f(w_t),
\qquad
w_{t+1}=w_t-\eta_t v_t.
\]
Common output is averaged iterate $\bar w=\frac{1}{T}\sum_{t=1}^T w_t$.

\subsection*{Variants}
\begin{itemize}
\item Decreasing step size, e.g. $\eta_t\propto 1/\sqrt t$
\item Strongly convex case often uses $\eta_t\propto 1/(\lambda t)$
\item Projected SGD enforces constraints by projection:
\[
w_{t+1}=\Pi_{\mathcal H}(w_t-\eta_t v_t).
\]
\end{itemize}

\subsection*{Convergence message from lecture}
For convex-Lipschitz bounded setup, SGD with suitable $T$ and $\eta$ reaches expected suboptimality $\le\varepsilon$ with sample complexity matching regularized ERM order.

\section{Perceptron + Feature Engineering Notes (\texttt{IE\_206\_Perceptron\_4.pdf})}
\subsection*{Halfspaces and affine score}
\[
h_{w,b}(x)=\langle w,x\rangle+b,
\qquad
\hat y=\operatorname{sign}(h_{w,b}(x))\in\{-1,+1\}.
\]

\subsection*{Perceptron update rule}
Given a mistake on $(x_i,y_i)$ where $y_i\langle w_t,x_i\rangle\le0$:
\[
w_{t+1}=w_t+y_ix_i.
\]
Repeat until no mistakes (if separable).

\subsection*{Convergence guarantee}
For separable data with
\[
R=\max_i\|x_i\|,
\quad
B=\min\{\|w\|:\ y_i\langle w,x_i\rangle\ge1\ \forall i\},
\]
lecture bound:
\[
T\le(RB)^2.
\]
Practical implication: feature scaling strongly affects convergence speed.

\subsection*{Normalization and centralization}
Standardization (z-score form):
\[
\phi(x)=\frac{x-\bar x}{\sigma}.
\]
Benefits: balanced feature influence, faster optimization, better numerical behavior.

\subsection*{Feature engineering and polynomial basis}
Map original input into richer features to make classes linearly separable:
\[
\phi(x_1,x_2)=(1,x_1,x_2,x_1^2,x_1x_2,x_2^2,\dots).
\]
XOR-style patterns become separable in transformed space.

\section{SVM Notes (\texttt{IE\_206\_SVM.pdf})}
\subsection*{Margin geometry}
Distance from point $x$ to hyperplane $(w,b)$:
\[
\operatorname{dist}(x,\{z:w^\top z+b=0\})=\frac{|w^\top x+b|}{\|w\|}.
\]
For canonical SVM constraints $y_i(w^\top x_i+b)\ge1$, margin width is
\[
\frac{2}{\|w\|}.
\]

\subsection*{Hard-margin SVM}
\[
\min_{w,b}\frac12\|w\|^2
\quad\text{s.t.}\quad
y_i(w^\top x_i+b)\ge1\ \forall i.
\]
Needs linear separability.

\subsection*{Soft-margin SVM}
Introduce slack $\xi_i\ge0$:
\[
\min_{w,b,\xi}\frac12\|w\|^2 + C\sum_i\xi_i
\quad\text{s.t.}\quad
y_i(w^\top x_i+b)\ge1-\xi_i.
\]
Equivalent regularized hinge-loss form:
\[
\min_{w,b}\lambda\|w\|^2+\frac1m\sum_i\max\{0,1-y_i(w^\top x_i+b)\}.
\]

\subsection*{Support vectors}
Points with active margin constraints determine the separator; in dual view, these are points with nonzero dual coefficients.

\subsection*{Duality and dual optimization}
For hard margin, dual (schematic):
\[
\max_{\alpha\ge0}\sum_i\alpha_i-\frac12\sum_{i,j}\alpha_i\alpha_j y_i y_j\langle x_i,x_j\rangle.
\]
Key implications:
\begin{itemize}
\item solution is linear combination of training points
\item optimization depends only on inner products (kernel entry point)
\item weak duality always; strong duality for convex feasible settings
\end{itemize}

\section{KNN + Bayes Notes (\texttt{IE\_206\_KNN.pdf})}
\subsection*{K-Nearest Neighbors}
Workflow:
\begin{enumerate}
\item choose $k$ and a distance metric
\item compute distance from query to all train points
\item pick $k$ nearest points
\item classify by majority vote (or average for regression)
\end{enumerate}

Common distances for $x,y\in\R^n$:
\[
\|x-y\|_2,\quad \|x-y\|_1,\quad \|x-y\|_p.
\]

Limitations from lecture:
\begin{itemize}
\item curse of dimensionality
\item prediction-time cost grows with dataset size
\item sensitivity to noisy data and scaling
\end{itemize}

\subsection*{Bayes classifier}
Bayes theorem:
\[
\PP(Y\mid X)=\frac{\PP(X\mid Y)\PP(Y)}{\PP(X)}.
\]
Optimal classifier:
\[
h^*(x)=\arg\max_y \PP(Y=y\mid X=x).
\]
Binary case:
\[
h^*(x)=\1\{\eta(x)>0.5\},\quad \eta(x)=\PP(Y=1\mid X=x).
\]
Bayes risk (minimum possible classification error):
\[
R^*=\E_X\left[1-\max_k\PP(Y=k\mid X)\right].
\]

\subsection*{Naive Bayes}
Conditional independence assumption:
\[
\PP(Y\mid X_1,\dots,X_n)\propto \PP(Y)\prod_{j=1}^n \PP(X_j\mid Y).
\]
Fast, scalable, often strong baseline (especially text tasks).

\subsection*{Discriminative vs generative}
\begin{itemize}
\item Discriminative: model $\PP(Y\mid X)$ or decision boundary directly (e.g., SVM, logistic)
\item Generative: model joint/conditional data generation (e.g., Naive Bayes)
\end{itemize}

\section{Multiclass Notes (\texttt{IE\_206\_MultiClass.pdf})}
\subsection*{Problem}
Predict one of $k$ classes for each input.

\subsection*{Reduction methods}
\textbf{One-vs-All (OvA):}
train $k$ binary models $h_i$; predict
\[
\hat y=\arg\max_i h_i(x).
\]

\textbf{All-Pairs (OvO):}
train one binary model per class pair $(i,j)$; aggregate pairwise scores/votes.

\subsection*{Pitfalls of reductions}
For imbalanced/low-mass classes, OvA can output trivial all-negative model for a class, causing severe multiclass failure even when a global multiclass linear separator exists.

\subsection*{Linear multiclass predictor}
Unified score model:
\[
\hat y=\arg\max_{y\in\{1,\dots,k\}} \langle w,\Psi(x,y)\rangle.
\]
Advantages: one jointly trained model; supports cost-sensitive losses $\Delta(y',y)$.

\subsection*{Multiclass SVM + SGD}
Structured hinge objective (lecture form):
\[
\min_w\lambda\|w\|^2+\frac1m\sum_{i=1}^m \max_{y'}\Big[\Delta(y',y_i)+\langle w,\Psi(x_i,y')-\Psi(x_i,y_i)\rangle\Big].
\]
Typical SGD step uses worst offending class $y^*$:
\[
w_{t+1}=w_t-\eta_t\big(\Psi(x,y^*)-\Psi(x,y)\big)
\]
(with regularization term as applicable).

\subsection*{Ranking intro}
Output is an ordering, not a single class. Mentioned losses: Kendall-tau pairwise disagreement and NDCG (top-heavy relevance quality).

\section{Decision Trees, Random Forest, Boosting Notes (\texttt{IE\_206\_RandoMForest.pdf})}
\subsection*{Decision trees}
Hierarchical split model with nodes, branches, leaves. Greedy top-down building is used because globally optimal tree search is computationally hard.

\subsection*{ID3-style construction}
Recursive rules:
\begin{itemize}
\item stop if pure node
\item stop if no features left; output majority label
\item otherwise split using feature with best gain and recurse
\end{itemize}

\subsection*{Gain measures}
For binary class proportion $a$:
\[
\text{Entropy }C(a)=-a\log a-(1-a)\log(1-a),
\]
\[
\text{Gini }C(a)=2a(1-a).
\]
Information gain = impurity before split minus weighted impurity after split.

\subsection*{Real-valued features via thresholds}
Convert to binary tests $\1[x_j<\theta]$ with candidate thresholds from sorted feature values. Efficient gain search can be implemented in about $O(dm\log m)$ (or better with incremental updates after sorting).

\subsection*{Random forest}
Core idea: reduce variance by averaging many decorrelated trees.
\begin{itemize}
\item bootstrap sample of training data per tree
\item random subset of features at each split
\item majority vote aggregation
\end{itemize}
Small feature subset size per split helps decorrelate trees.

\subsection*{Boosting and decision stumps}
Decision stump: one-split tree (weak learner).
Boosting trains weak learners sequentially with reweighting to focus on hard examples.

\subsection*{ERM for decision stumps}
Given weighted sample distribution $D$, choose feature-threshold pair minimizing weighted classification error. Sorting by each feature and sweeping thresholds yields efficient implementation.

\subsection*{AdaBoost}
At round $t$:
\begin{itemize}
\item train weak learner $h_t$ on distribution $D_t$
\item weighted error $\varepsilon_t$
\item learner weight $\alpha_t=\frac12\log\frac{1-\varepsilon_t}{\varepsilon_t}$
\item upweight misclassified examples for next round
\end{itemize}
Final model:
\[
H(x)=\operatorname{sign}\left(\sum_{t=1}^T \alpha_t h_t(x)\right).
\]
Training error bound (lecture theorem style): decays exponentially with rounds when each weak learner beats random guessing ($\varepsilon_t<1/2$).

\section{Practice Sheet Notes (\texttt{IE206\_PracticeQs\_1.pdf}, \texttt{IE206\_PracticeQs\_2.pdf})}
\subsection*{PracticeQs\_1 themes}
\begin{itemize}
\item \textbf{Cost-sensitive decision making}: choose model by total expected penalty, not by accuracy alone.
\item \textbf{Imbalanced fraud setting}: compare accuracy vs precision/recall/F1 and pick metric matching business objective.
\item \textbf{LOOCV for linear regression}: compute leave-one-out models and compare fit quality ($R^2$/MSE).
\item \textbf{Hyperparameter tuning}: choose setting with best validation performance (generalization), avoid overfitting settings with low train but poor validation error.
\item \textbf{Sherman--Morrison identity}: rank-one inverse update useful for efficient LOOCV derivation.
\item \textbf{Fast LOOCV formula via hat matrix}:
\[
\hat y_i^{(-i)}=\frac{\hat y_i-h_{ii}y_i}{1-h_{ii}},
\qquad
e_i^{(-i)}=\frac{e_i}{1-h_{ii}},
\]
so
\[
\text{CV}(n)=\frac1n\sum_{i=1}^n\left(\frac{y_i-\hat y_i}{1-h_{ii}}\right)^2.
\]
\item \textbf{Clipped/rectified regression loss}:
\[
\ell(w;x,y)=\big|y-\max(0,w^\top x)\big|.
\]
Convex in $w$ as composition of convex functions with affine map plus absolute value; non-differentiable at hinge boundaries where $w^\top x=0$ and/or where the outer absolute value argument is zero.
\end{itemize}

\subsection*{PracticeQs\_2 themes}
\begin{itemize}
\item \textbf{Bayes vs 1-NN asymptotics}: pointwise error bound
\[
\varepsilon_{NN}(x)\le2\varepsilon_{Bayes}(x).
\]
\item \textbf{Feature scaling in KNN}: nearest neighbor can flip after normalization due to scale dominance in raw Euclidean distance.
\item \textbf{Hard-margin SVM geometry}: separability proof, optimal separator, support vector identification, effect of removing non-support vectors.
\item \textbf{Soft-margin extremes}: behavior for $C=0$ (regularization dominates) and $C\to\infty$ (hard-margin-like behavior if separable).
\item \textbf{KKT interpretation in soft SVM}: from
\[
\alpha_i(1-y_if_i-\xi_i)=0,\quad \beta_i\xi_i=0,\quad \alpha_i+\beta_i=C,
\]
if $\alpha_i=C$ then $\beta_i=0$; point is inside margin or misclassified (cannot be strictly outside margin with inactive slack and $\alpha_i=C$ simultaneously).
\end{itemize}

\section{Quick Exam-Revision Checklist}
\begin{enumerate}
\item Derive OLS and ridge closed forms, and know when matrix inverse fails.
\item Compute and interpret $R^2$, SST/SSM/SSE, and correlation.
\item Know robust losses and when to prefer L1/Huber over L2.
\item Use precision/recall/F1 correctly for imbalanced data.
\item Implement CV logic and LOOCV hat-matrix shortcut.
\item Remember perceptron mistake update and convergence condition.
\item Derive hard/soft SVM primal forms, hinge loss form, and dual interpretation.
\item Explain KNN distance scaling and curse of dimensionality.
\item Distinguish Bayes optimal, Naive Bayes, discriminative vs generative.
\item Compare OvA/OvO vs joint multiclass methods.
\item Explain random forest variance reduction mechanism.
\item State AdaBoost update logic and why training error can drop exponentially.
\item Explain SGD vs GD and when projection/subgradient is required.
\end{enumerate}

\end{document}
